% Hong Tang -- Curriculum Vitae
% Source maintained for the homepage CV page.

\documentclass[11pt,a4paper]{article}
\usepackage[margin=1.8cm]{geometry}
\usepackage[hidelinks]{hyperref}
\usepackage{enumitem}
\setlist[itemize]{leftmargin=1.2em,itemsep=2pt,topsep=2pt}
\pagenumbering{gobble}

\begin{document}

\begin{center}
{\Large\textbf{Hong Tang}}\\
Research Assistant Professor, Department of Building Environment and Energy Engineering\\
The Hong Kong Polytechnic University\\
\href{mailto:honghong.tang@polyu.edu.hk}{honghong.tang@polyu.edu.hk} \quad
\href{https://scholar.google.com.hk/citations?user=nXW7KrwAAAAJ&hl=en}{Google Scholar} \quad
\href{https://www.researchgate.net/profile/Hong-Tang-17}{ResearchGate}
\end{center}

\section*{Research Interests}
Grid-interactive buildings; building energy flexibility and demand response; AI and machine learning for energy systems; data-center precision cooling and computing-power/electricity coordination; integrated energy systems and virtual power plants.

\section*{Education}
\begin{itemize}
\item PhD, Building Environment and Energy Engineering, The Hong Kong Polytechnic University, 2018--2023.
\item BEng, Building Environment and Energy Application Engineering, Tsinghua University, 2014--2018.
\end{itemize}

\section*{Appointments}
\begin{itemize}
\item Research Assistant Professor, The Hong Kong Polytechnic University, 2025--present.
\item Postdoctoral Fellow, The Hong Kong Polytechnic University, 2023--2025.
\item Research Assistant, The Hong Kong Polytechnic University, 2022--2023.
\end{itemize}

\section*{Selected Publications}
\begin{itemize}
\item Hong Tang et al. Multi-Timescale Rolling Optimization Coordinating Heterogeneous EV and Building Load Flexibility in Community Energy Systems with NGBoost-Based Renewable Forecast. \emph{Applied Energy}, accepted, 2026.
\item Hong Tang et al. Risk-aware optimal dispatch of resource aggregators integrating NGBoost-based probabilistic renewable prediction and bi-level building flexibility engagement. \emph{Engineering}, 53:76--89, 2025.
\item Hong Tang et al. Unlocking Building Energy Flexibility: Multi-Layered Pricing Framework Design from Distribution-Level to P2P Markets. \emph{The Innovation Energy Use}, 1(2):100019, 2025.
\item Hong Tang and Shengwei Wang. Game-theoretic optimization of demand-side flexibility engagement considering the perspectives of different stakeholders and multiple flexibility services. \emph{Applied Energy}, 332:120550, 2023.
\item Hong Tang, Shengwei Wang and Hangxin Li. Flexibility categorization, sources, capabilities and technologies for energy-flexible and grid-responsive buildings. \emph{Energy}, 219:119598, 2021.
\end{itemize}

\section*{Service and Awards}
\begin{itemize}
\item Young Editorial Board Member, \emph{The Innovation Energy Use}.
\item Guest Editor for \emph{Energy}, \emph{Building Simulation}, \emph{Buildings}, \emph{Architecture}, \emph{Journal of Physics: Conference Series}, and \emph{Journal of Building Design and Environment}.
\item Best Presentation Award, The 11th Asia Conference on Power and Electrical Engineering, 2026.
\item Best Poster Award, The 4th Asia Conference of International Building Performance Simulation Association, 2018.
\end{itemize}

\vfill
\noindent Last updated: May 27, 2026

\end{document}
